\documentclass[12pt,a4paper]{article}

% --- CONFIGURACION PARA OVERLEAF (UTF8) ---
\usepackage[T1]{fontenc}
\usepackage[utf8]{inputenc}
\usepackage[spanish]{babel}
\usepackage{csquotes}
\usepackage{graphicx}
\usepackage{amsmath}
\usepackage{amssymb}
\usepackage{booktabs}
\usepackage{hyperref}
\usepackage{float}
\usepackage{listings}
\usepackage{xcolor}
\usepackage{geometry}
\usepackage{fancyhdr}
\usepackage{setspace}

% Configuracion de margenes
\geometry{left=2.5cm, right=2.5cm, top=2.5cm, bottom=2.5cm}

% Configuracion de encabezado y pie
\pagestyle{fancy}
\fancyhf{}
\rhead{Informe Predictivo}
\lhead{Prediccion de Precios de Viviendas}
\cfoot{\thepage}

% Configuracion de listings
\lstset{
    language=Python,
    basicstyle=\ttfamily\small,
    keywordstyle=\color{blue},
    commentstyle=\color{gray},
    stringstyle=\color{red},
    breaklines=true,
    showstringspaces=false,
    backgroundcolor=\color{gray!10}
}

\title{\textbf{Informe Tecnico: Analisis y Prediccion de Precios de Viviendas}}
\author{Analisis Computacional de Datos Inmobiliarios}
\date{Febrero, 2026}

\begin{document}

\maketitle

\begin{abstract}
Este informe presenta un analisis exhaustivo de un modelo de prediccion de precios de viviendas basado en aprendizaje automatico. El proyecto procesa y analiza un conjunto de datos que contiene aproximadamente 1 millon de registros inmobiliarios, utilizando caracteristicas como tamano de la propiedad, numero de habitaciones, antiguedad y ubicacion geografica. Se evaluaran las ventajas y desventajas del enfoque implementado, asi como sus implicaciones practicas.
\end{abstract}

\tableofcontents
\newpage

% ========================================
\section{Introduccion}
% ========================================

La prediccion de precios en el mercado inmobiliario es un problema de regresion clasico en la ciencia de datos. Con la creciente disponibilidad de datos sobre propiedades y transacciones, los modelos de aprendizaje automatico ofrecen oportunidades sin precedentes para estimar valores de manera precisa y automatizada.

Este proyecto aborda la construccion de un modelo predictivo escalable capaz de procesar y analizar un volumen significativo de datos (aproximadamente 1 millon de filas) para generar predicciones confiables de precios de viviendas.

% ========================================
\section{Descripcion del Problema}
% ========================================

\subsection{Objetivo}
Desarrollar un modelo de regresion que prediga el precio de una vivienda dado un conjunto de caracteristicas independientes.

\subsection{Variables del Modelo}
\begin{table}[H]
\centering
\begin{tabular}{|l|l|l|}
\hline
\textbf{Variable} & \textbf{Tipo} & \textbf{Descripcion} \\
\hline
Tamano & Numerica (continua) & Area de la propiedad en m2 \\
Habitaciones & Numerica (discreta) & Numero de cuartos \\
Antiguedad & Numerica (continua) & Anos desde su construccion \\
Zona & Categorica & Ubicacion geografica (A, B, C, etc.) \\
Precio & Numerica (continua) & Variable dependiente a predecir \\
\hline
\end{tabular}
\caption{Variables del conjunto de datos}
\end{table}

\subsection{Escala del Proyecto}
El dataset contiene aproximadamente \textbf{1,000,000 de registros}, lo que presenta desafios significativos en terminos de:
\begin{itemize}
    \item Almacenamiento y procesamiento de datos
    \item Tiempo de entrenamiento del modelo
    \item Memoria computacional requerida
    \item Optimizacion y escalabilidad
\end{itemize}

% ========================================
\section{Metodologia}
% ========================================

\subsection{Preprocesamiento de Datos}
El preprocesamiento incluye las siguientes etapas:

\begin{enumerate}
    \item \textbf{Limpieza de datos}: Eliminacion de valores nulos y duplicados
    \item \textbf{Codificacion de variables categoricas}: Conversion de variables de zona a formato numerico
    \item \textbf{Normalizacion}: Escalado de caracteristicas numericas
    \item \textbf{Division del conjunto}: Separacion en datos de entrenamiento (80\%) y prueba (20\%)
\end{enumerate}

\subsection{Algoritmos Considerados}
Se contemplan los siguientes enfoques para modelado:

\begin{table}[H]
\centering
\begin{tabular}{|l|l|l|}
\hline
\textbf{Algoritmo} & \textbf{Ventaja} & \textbf{Limitacion} \\
\hline
Regresion Lineal & Simplicidad, interpretabilidad & Linealidad asumida \\
Random Forest & Robustez, manejo de no-linealidad & Tiempo de entrenamiento \\
Gradient Boosting & Precision alta & Complejidad del modelo \\
Neural Networks & Flexibilidad extrema & Riesgo de sobreajuste \\
\hline
\end{tabular}
\caption{Comparativa de algoritmos de regresion}
\end{table}

\subsection{Metricas de Evaluacion}
Los modelos seran evaluados mediante:

\begin{equation}
\text{MAE} = \frac{1}{n}\sum_{i=1}^{n}|y_i - \hat{y}_i|
\end{equation}

\begin{equation}
\text{RMSE} = \sqrt{\frac{1}{n}\sum_{i=1}^{n}(y_i - \hat{y}_i)^2}
\end{equation}

\begin{equation}
R^2 = 1 - \frac{\sum_{i=1}^{n}(y_i - \hat{y}_i)^2}{\sum_{i=1}^{n}(y_i - \bar{y})^2}
\end{equation}

Donde $y_i$ son valores reales, $\hat{y}_i$ son predicciones y $\bar{y}$ es la media.

% ========================================
\section{Ventajas del Modelo}
% ========================================

\subsection{Escalabilidad}
\begin{itemize}
    \item El modelo puede procesarse con frameworks como Spark o Dask para distribuir la carga computacional
    \item Permite integracion en pipelines de produccion
    \item Manejo eficiente de datasets masivos mediante chunking
\end{itemize}

\subsection{Precision Predictiva}
\begin{itemize}
    \item Multiples caracteristicas permiten capturar la complejidad del mercado inmobiliario
    \item Con 1M de muestras, el modelo tiene suficientes datos para generalizar bien
    \item Reduccion del error mediante tecnicas de validacion cruzada
\end{itemize}

\subsection{Interpretabilidad}
\begin{itemize}
    \item Las caracteristicas utilizadas son directamente interpretables para stakeholders
    \item Posibilidad de analizar importancia de caracteristicas
    \item Explicacion clara del impacto de cada variable en el precio
\end{itemize}

\subsection{Flexibilidad}
\begin{itemize}
    \item El modelo puede adaptarse a diferentes mercados geograficos
    \item Facil incorporacion de nuevas caracteristicas (energia, seguridad, amenidades)
    \item Reentrenamiento periodico con datos actualizados
\end{itemize}

% ========================================
\section{Desventajas del Modelo}
% ========================================

\subsection{Limitaciones Tecnicas}
\begin{itemize}
    \item \textbf{Overhead computacional}: El entrenamiento de modelos complejos con 1M de registros requiere infraestructura costosa
    \item \textbf{Tiempo de ejecucion}: Modelos como Gradient Boosting pueden tardar horas en entrenar
    \item \textbf{Dependencia de hardware}: Requiere suficiente RAM para cargar datasets masivos
\end{itemize}

\subsection{Limitaciones de Datos}
\begin{itemize}
    \item \textbf{Calidad del dato}: Errores en la recopilacion de datos se amplifican con 1M de registros
    \item \textbf{Sesgos historicos}: El modelo aprende sesgos del mercado pasado
    \item \textbf{Volatilidad del mercado}: Cambios economicos rapidos pueden invalidar predicciones
    \item \textbf{Informacion faltante}: Caracteristicas adicionales relevantes no disponibles
\end{itemize}

\subsection{Riesgos de Generalización}
\begin{itemize}
    \item \textbf{Sobreajuste}: Con dataset tan grande, el modelo podria captar ruido
    \item \textbf{Distribucion sesgada}: Si los datos no representan bien todos los segmentos del mercado
    \item \textbf{Cambios temporales}: Patrones historicos pueden no aplicar al presente
\end{itemize}

\subsection{Factores Externos No Capturados}
\begin{itemize}
    \item Cambios en politicas de vivienda
    \item Fluctuaciones economicas macro
    \item Eventos locales (desarrollos, infraestructura)
    \item Preferencias subjetivas del comprador
    \item Caracteristicas intangibles de la propiedad
\end{itemize}

% ========================================
\section{Analisis Comparativo de Resultados}
% ========================================

\subsection{Escenario Optimista}
Si el modelo logra un $R^2 \geq 0.85$:

\begin{table}[H]
\centering
\begin{tabular}{|l|l|}
\hline
\textbf{Aspecto} & \textbf{Implicacion} \\
\hline
Precision & Modelo altamente confiable \\
Aplicacion & Viable para decisiones financieras \\
Valor Comercial & Integracion en plataformas inmobiliarias \\
\hline
\end{tabular}
\end{table}

\subsection{Escenario Realista}
Si el modelo logra un $0.70 \leq R^2 < 0.85$:

\begin{table}[H]
\centering
\begin{tabular}{|l|l|}
\hline
\textbf{Aspecto} & \textbf{Implicacion} \\
\hline
Precision & Modelo util con margenes de error \\
Aplicacion & Apoyo a decisiones, no determinante \\
Valor Comercial & Herramienta complementaria \\
\hline
\end{tabular}
\end{table}

\subsection{Escenario Pesimista}
Si el modelo logra un $R^2 < 0.70$:

\begin{table}[H]
\centering
\begin{tabular}{|l|l|}
\hline
\textbf{Aspecto} & \textbf{Implicacion} \\
\hline
Precision & Modelo con limitaciones significativas \\
Aplicacion & Requiere refinamiento de caracteristicas \\
Valor Comercial & Necesidad de recolectar mas datos \\
\hline
\end{tabular}
\end{table}

% ========================================
\section{Recomendaciones}
% ========================================

\subsection{Para Mejora del Modelo}
\begin{enumerate}
    \item Incorporar mas caracteristicas relevantes (proximidad a transporte, servicios, seguridad)
    \item Implementar ingenieria de caracteristicas (interacciones entre variables)
    \item Utilizar validacion cruzada estratificada
    \item Considerar modelos de ensemble para mejorar robustez
\end{enumerate}

\subsection{Para Implementacion en Produccion}
\begin{enumerate}
    \item Establecer pipelines de datos automatizados
    \item Implementar monitoreo de desempeno del modelo
    \item Crear alertas para cambios significativos en el mercado
    \item Planificar reentrenamiento periodico (mensual/trimestral)
    \item Documentar supuestos del modelo y limitaciones
\end{enumerate}

\subsection{Para Gestion de Riesgo}
\begin{enumerate}
    \item Proporcionar intervalos de confianza junto a predicciones
    \item Mantener auditoria de decisiones tomadas basadas en el modelo
    \item Establecer limites de confianza para aplicar el modelo
    \item Supervision humana para casos fuera de rango normal
\end{enumerate}

% ========================================
\section{Conclusiones}
% ========================================

El modelo de prediccion de precios de viviendas basado en un dataset de 1 millon de registros representa una oportunidad significativa para automatizar y mejorar la precision en valuación inmobiliaria. Sin embargo, sus limitaciones tecnicas y el factor humano en el mercado inmobiliario requieren que sea visto como una herramienta de apoyo, no como verdad absoluta.

\subsection{Puntos Clave}
\begin{itemize}
    \item \textbf{Ventaja principal}: Escalabilidad y capacidad de procesamiento masivo de datos
    \item \textbf{Desventaja principal}: Volatilidad del mercado y factores no cuantificables
    \item \textbf{Aplicación realista}: Herramienta de apoyo en decisiones inmobiliarias
    \item \textbf{Camino futuro}: Integracion con variables economicas y analisis temporal
\end{itemize}

\subsection{Proximas Fases}
\begin{enumerate}
    \item Implementar y entrenar el modelo
    \item Evaluar metricas contra benchmarks de la industria
    \item Validar con datos externos y recientes
    \item Desarrollar interfaz de usuario para predicciones
    \item Establecer proceso de actualizacion del modelo
\end{enumerate}

% ========================================
\section{Apendice: Codigo Base}
% ========================================

\subsection{Carga y Preprocesamiento}
\begin{lstlisting}
import pandas as pd
import numpy as np
from sklearn.model_selection import train_test_split
from sklearn.preprocessing import StandardScaler

# Carga de datos
df = pd.read_csv('casas.csv')

# Analisis exploratorio
print(df.head())
print(df.info())
print(df.describe())

# Separacion de caracteristicas y target
X = df[['tamano', 'habitaciones', 'antiguedad']]
y = df['precio']

# Codificacion de variables categoricas
X = pd.get_dummies(X, columns=['zona'])

# Normalización
scaler = StandardScaler()
X_scaled = scaler.fit_transform(X)

# Division entrenamiento/prueba
X_train, X_test, y_train, y_test = train_test_split(
    X_scaled, y, test_size=0.2, random_state=42
)
\end{lstlisting}

\subsection{Entrenamiento de Modelo}
\begin{lstlisting}
from sklearn.ensemble import RandomForestRegressor
from sklearn.metrics import mean_absolute_error, r2_score

# Entrenamiento
model = RandomForestRegressor(n_estimators=100, 
                               random_state=42, 
                               n_jobs=-1)
model.fit(X_train, y_train)

# Predicciones
y_pred = model.predict(X_test)

# Evaluacion
mae = mean_absolute_error(y_test, y_pred)
r2 = r2_score(y_test, y_pred)

print(f'MAE: ${mae:.2f}')
print(f'R2 Score: {r2:.4f}')
\end{lstlisting}

\end{document}
